\section{Introduction}

Le problème du signe de la forme de Weil réunit deux exigences :
conserver toutes ses contributions arithmétiques et contrôler la forme
dans des espaces complets de fonctions. Le critère global relie
sa positivité à l’hypothèse de Riemann. Une construction
spectrale qui vise ce critère doit donc spécifier
le produit, les poids et les domaines qui déterminent la forme,
et démontrer comment ses termes croisés se transmettent. Cet
article développe cette relation à partir d’une arithmétique de chemins
avec mémoire et obtient des résultats exacts de reconstruction,
de transport et de coercivité.

La principale contribution fonctionnelle est une réduction de Schur pour chaque fenêtre fixée. Une partition suffisamment fine sépare les moyennes cellulaires d'un espace infini de détails de moyenne nulle dont la forme est coercive. L'élimination exacte de ces détails conserve tout leur couplage avec les moyennes et produit une matrice finie ayant le même indice négatif que la forme complète. Le résultat donne un critère exact, et non un remplacement de l'espace des fonctions par des échantillons. Son application dans la fenêtre $(-1/18,1/18)$ démontre une borne inférieure de $\frac12\|f\|_2^2$ pour tout le domaine conjoint de ses lecteurs. Un transport réversible prolonge cette borne à des collections complètes de fonctions tests à queues exponentielles.

\subsection*{Construction arithmétique et structure conservée}

Le point de départ est un noyau formel mathématique appelé
holographie modulaire triadique, ci-après HMT. Son objet
\emph{All Possible Paths} (APP) réunit des chemins et des évaluations
arithmétiques sur le support $X_9=(\mathbb Z/9\mathbb Z)^2$.
Les générateurs cardinaux $\{(1,0),(-1,0),(0,1),(0,-1)\}$
définissent le graphe de Cayley de son groupe additif, un réseau
cyclique bidimensionnel à réalisation toroïdale. Sur les mêmes
positions agissent les évaluations somme et produit, chacune
avec son résidu et son quotient. Le support spécifie où une
route se déplace ; les feuillets arithmétiques spécifient ce qui est
calculé et conservé pendant ce déplacement.

TRIT est la signature ternaire orientée $\{-1,0,+1\}$, obtenue
au moyen de représentants équilibrés modulo $3$. Elle distingue les
types hyperbolique, parabolique et elliptique dans sa réalisation
quadratique. Le \emph{Topological Phase Kernel}, ci-après
TPK, compose sélection de phase, transport et actualisation de
la mémoire. Son état conserve positions, feuillets, orientation,
quotients, préfixes et frontière. Le prolongement nonadique transporte
cette information entre profondeurs : le retour de phase fait avancer
l’histoire au lieu de la réinitialiser.

La structure discrète conjointe du continu se développe sur ces mêmes états et leurs compatibilités. Ses constructions corrélatives conservent restrictions, incidences, formes et mémoire dans un unique système de raffinement. La lecture archimédienne obtient des valeurs à partir de cylindres compatibles et la lecture de frontière retient les incidences utilisées par les réalisations ultérieures. Le présent problème arithmétique et spectral conserve cette origine commune, avec les opérations et les domaines qui interviennent effectivement dans chaque résultat.

La relation entre centre et frontière fournit une première instance de reconstruction. Dans la coordinatisation multiaffine d'APP, la frontière relevée détermine l'intérieur ; les valeurs résiduelles sont complétées par les quotients qui distinguent leurs valeurs entières. Les signatures de continuation exigent de conserver les feuillets, et le calendrier des vacances conserve l'écart entre capacité et profondeur. La section~\ref{sec:centro-frontera} situe dans cette construction le centre $(5,5)$, la signature régionale $555555$, les lecteurs $72,27,77,22$ et les intervalles $21/22$ et $6/7$. Leur fonction spécifie l'information transportable avant toute identification spectrale.

\subsection*{Du produit à la fonction complétée}

Le produit natif compose les deux feuillets d’APP au moyen d’un
transducteur de chiffres et de quotients. Son invariant d’évaluation
prouve l’exactitude de l’opération sur les mots nonadiques,
et l’unicité de la forme normale permet de reconnaître ensuite la
multiplication entière. Les fibres du produit déterminent ses
ouvertures ; l’absence d’une ouverture non triviale sélectionne
les irréductibles. Le théorème~\ref{rz-num-01-6} identifie leurs
évaluations aux nombres premiers positifs.

La factorisation et les horloges arithmétiques décrivent des propriétés distinctes de ces objets. La première détermine leurs composantes multiplicatives ; les secondes enregistrent des retours d'orbites résiduelles. La représentation par occupations de modes conserve la multiplicité de chaque irréductible et conduit au produit d'Euler. La trace sur toutes les occupations et la trace sur les modes individuels sont reliées par les moments premier--puissance de la proposition~\ref{rz-num-02-3}.

Le raffinement cyclique fournit ensuite une trace de
chaleur dont la limite est la série thêta. La transformation de Mellin
produit $\pi^{-s/2}\Gamma(s/2)Z(s)$ ; l’inversion de l’échelle
détermine la contribution polaire et l’équation fonctionnelle de la
forme complétée. Les propositions~\ref{rz-num-03-1}
et~\ref{rz-num-03-2} maintiennent la normalisation et la convergence
de cette étape. Les développements des parties finies et des lecteurs
résiduels étendent leurs évaluations avec des restes contrôlés.
La forme de Weil conserve conjointement les termes premier,
gamma et polaire qui proviennent de cette composition.

\subsection*{Double cercle, reconstruction et transport}

Le double cercle distingue la coordonnée spectrale de la mémoire
des tours. La décomposition en neuf phases produit un bilan
exact de coefficient $8/81$ ; les modes du bloc central
déterminent le rapport $q=5/9$. Les moments de mémoire
$q^{|n|}$ possèdent une factorisation triangulaire positive en toute
dimension. Le registre des pertes conserve en outre le terme
terminal et les produits croisés, de sorte que la contraction
visible peut s’accompagner d’une conservation isométrique de
l’histoire enregistrée.

La réalisation analytique construit, sur des fonctions tests à décroissance exponentielle, un opérateur réversible $C_a$, avec $a>b>1/2$. Son action intégrale conserve simultanément les corrélations et les produits polaires croisés. L'identité de la sous-section~\ref{tm-add:invariancia-weil} démontre
\[
\mathscr W(C_af,C_ag)=\mathscr W(f,g).
\]
La marge $a>1/2$ appartient au domaine de ce lecteur analytique
et garantit la convergence de ses contributions polaires.
Le cercle régulier de mémoire et une représentation qui reçoit
une limite atomique conservent leurs domaines respectifs.

La queue gamma contient un lecteur inverse explicite : il retrouve
la fonction test et permet de reconstruire la colonne négative
à partir de la positive. Les équations~\eqref{gamma-add-eq-12}
et~\eqref{gamma-add-eq-14} décrivent son action, et
\eqref{gamma-add-eq-16} conserve la différence complète de
formes. Cette reconstruction donne l’opérateur à étudier ;
son caractère borné et sa contractivité sont des propriétés différentes,
dont les estimations sont développées ci-après.

\subsection*{Coercivité et critère de prolongement global}

Le théorème~\ref{schur-add-schur-exacto} démontre qu'une fois obtenue la borne positive du bloc de détails $D$, la positivité de la forme sur $V_R$ équivaut à celle de $S=A-B^*D^{-1}B$. La correction conserve l'effet de tous les détails et de leurs termes croisés avec les moyennes. Le théorème~\ref{schur-add-residual} encadre cette matrice au moyen de résidus et d'une borne opératorielle de l'erreur. Dans la fenêtre de longueur $1/9$, les estimations $D>I$, $A>97/100$ et $\|B\|^2<1/5$ conduisent au théorème~\ref{schur-add:coercividad-ventana}. La fenêtre locale satisfait $1/9<\log2$ ; ses collections transportées peuvent avoir des queues et reçoivent alors la somme première complète au moyen de l'invariance démontrée.

Le critère global exige également les termes croisés entre collections distinctes et la couverture de tout le domaine des fonctions tests. La formulation précise se trouve dans les sous-sections~\ref{tm-add-m6-2}, \ref{tm-add-m7} et~\ref{gamma-add-sub-375} : identifier les moments arithmétiques complets à un Gram positif équivaut à conserver simultanément l'énergie et le transport, avec leurs domaines et leurs limites. Cette condition demeure une hypothèse des résultats globaux ; la coercivité démontrée appartient aux domaines explicites précédents. L'apport de l'article consiste à unir la génération arithmétique à des identités de reconstruction et à des estimations fonctionnelles suffisamment précises pour établir cette portée et sa condition exacte de prolongement.
\clearpage
